\documentclass[11pt]{article}

\usepackage[margin=1in]{geometry}
\usepackage{amsmath,amssymb,amsthm}
\usepackage[T1]{fontenc}
\usepackage[utf8]{inputenc}
\usepackage{lmodern}

\title{Billingsley's Probability Measure Axiom and Collective Exhaustion}
\author{}
\date{\today}

\begin{document}
\maketitle

\section*{Classical starting point}

In Billingsley's formulation, a probability measure is a set function
$P$ on a field $\mathcal F$ of subsets of $\Omega$ satisfying:
\begin{enumerate}
  \item[(i)] $P(A) \ge 0$ for every $A \in \mathcal F$;
  \item[(ii)] $P(\Omega)=1$;
  \item[(iii)] if $A_1,A_2,\ldots \in \mathcal F$ are pairwise disjoint and
  $\bigcup_n A_n \in \mathcal F$, then
  \[
    P\!\left(\bigcup_{n=1}^{\infty} A_n\right)
    =
    \sum_{n=1}^{\infty} P(A_n).
  \]
\end{enumerate}

For a finitely additive probability, this third axiom is equivalent to
continuity at the empty set:
\[
  E_1 \supseteq E_2 \supseteq \cdots,
  \qquad
  \bigcap_n E_n = \emptyset
  \quad\Longrightarrow\quad
  P(E_n) \to 0.
\]
Thus the analytic content of countable additivity is that mass cannot
persist along a descending sequence of events that has vanished.

\section*{Paper I starting point}

Paper I begins one step earlier.  Instead of assuming a single
countably additive probability measure, it starts with a directed query
system: Boolean algebras $\mathcal E_i$ connected by refinement maps
$\pi_{ij}$, together with compatible finitely additive charges
$\{\ell_i\}$.

The structural input is a query system consisting of
\[
  \bigl(\mathcal I,\le\bigr),
  \qquad
  (O_i,\mathcal E_i)_{i\in\mathcal I},
  \qquad
  \pi_{ij}:O_j\to O_i\quad(i\le j).
\]
These data are required to satisfy the standing refinement identities
\[
  \pi_{ii}=\mathrm{id},\qquad
  \pi_{ik}=\pi_{ij}\circ\pi_{jk}\quad(i\le j\le k),
\]
and the relevant directedness hypotheses used by the chosen extension route.
No separate underlying sample space is needed at this stage.  When a realised
space is required, Paper I may take the canonical instantiation
\[
  \Omega
  =
  \varprojlim_{i\in\mathcal I}O_i
  =
  \left\{
    (\omega_i)_i\in\prod_i O_i:
    \pi_{ij}(\omega_j)=\omega_i\ \text{for }i\le j
  \right\},
\]
with evaluation maps $\mathrm{eval}_i(\omega)=\omega_i$.  Thus
\[
  \mathrm{eval}_i=\pi_{ij}\circ\mathrm{eval}_j\quad(i\le j)
\]
holds by construction.  The additional evaluation and separation hypotheses
used in the descent to a realised observable measure are:
\[
  \mathrm{eval}_i:\Omega\to O_i\ \text{is surjective},
\]
and, for the Stone route, discriminability:
\[
  \omega\ne\omega'
  \quad\Longrightarrow\quad
  \exists i\in\mathcal I,\ \exists A\in\mathcal E_i
  \ \text{such that}\
  \mathrm{eval}_i(\omega)\in A,\quad
  \mathrm{eval}_i(\omega')\notin A.
\]
The observable cylinder algebra is
\[
  \mathcal C
  =
  \{\mathrm{eval}_i^{-1}(A):i\in\mathcal I,\ A\in\mathcal E_i\}.
\]

The observable-data analogue of Billingsley's definition is not yet a
probability measure on a $\sigma$-algebra.  It is an admissible coherent
probability datum: a family $\{\ell_i\}$ satisfying the following
conditions.
\begin{enumerate}
  \item[(i)] Local finite probability: each $\ell_i$ is a normalised finitely
  additive charge on $\mathcal E_i$, i.e.
  \[
    \ell_i(A) \ge 0 \quad (A \in \mathcal E_i),
    \qquad
    \ell_i(O_i)=1,
  \]
  and for disjoint $A,B \in \mathcal E_i$,
  \[
    \ell_i(A \cup B)=\ell_i(A)+\ell_i(B);
  \]

  \item[(ii)] Compatibility under refinement:
  \[
    \ell_j(\pi_{ij}^{-1}A)=\ell_i(A)
    \qquad
    (i \le j,\; A \in \mathcal E_i);
  \]

  \item[(iii)] Collective exhaustion (CE): if $i \in \mathcal I$ and
  $E_1,E_2,\ldots \in \mathcal E_i$ satisfy
  \[
    E_1 \supseteq E_2 \supseteq \cdots,
    \qquad
    \bigcap_{n=1}^{\infty} E_n = \emptyset,
  \]
  then there exists $j \ge i$ such that
  \[
    \ell_j\!\left(\pi_{ij}^{-1}(E_n)\right) \longrightarrow 0.
  \]
\end{enumerate}

Here local finite probability expands into positivity, normalization, and finite
additivity at each level.  Compatibility is the new query-system coherence
condition.  CE replaces direct assumption of countable additivity at the
observable-data level: no compatible charge family may assign persistent mass to
finite observable conditions that collectively vanish under refinement.

These three conditions are the minimal measure-side hypotheses.  Together with
the structural query-system hypotheses above, they imply
\[
  \begin{aligned}
  &\exists!\,P:\sigma(\mathcal C)\to[0,1],\\
  &P\ \text{is countably additive},\\
  &(\mathrm{eval}_i)_\#P=\ell_i
    \qquad (i\in\mathcal I).
  \end{aligned}
\]

Thus Paper I does not alter the target notion of probability measure.  The
delivered object is an ordinary countably additive probability measure on
$(\Omega,\sigma(\mathcal C))$.  The change is in the input: countable
additivity is not assumed of a global $P$ at the outset; it is recovered from
local finite probability, refinement compatibility, and CE.

\end{document}
